\section{Замены координат в гамильтоновых системах}
\subsection{Производящая функция симплектического диффеоморфизма.}
Пусть есть гамильтонова система с гамильтонианом $H$ и каноническими координатами $\bt{x} = (\bt{q}, \bt{p})$.
Канонические уравнения записываются как
\[
    \dot{\bt{x}} = J\partial_{\bt{x}} H.
\]
Рассмотрим некоторый диффеоморфизм $\phi: (\bt{q}, \bt{p}) \to (\tilde{\bt{q}}, \tilde{\bt{p}})$.
Не всякий диффеоморфизм сохраняет гамильтоновость системы.
В этом параграфе мы попробуем определить и классифицировать преобразования, уважаюшие гамильтонову структуру динамической системы и посмотреть, что с ними (и как) можно сделать.
Для начала докажем маленькую техническую лемму, которая позволяет определять симплектичность матрицы:
\begin{lemma}
    Для произвольной матрицы $M$ (для которой имеют смысл написанные утверждения) следующие утверждения эквивалентны:
    \begin{itemize}
        \item Матрица $M$ -- симплектическая, то есть $M^T J M = J$.
        \item Представляя матрицу $M$ в виде блочной (из 4-х блоков $m \times m$) как
        \[
            M = \begin{pmatrix*} A_{11} & A_{12} \\ A_{21} & A_{22}
            \end{pmatrix*}
        \]
        Получаем следующие требования на блоки:
        \[
            \begin{cases}
                A_{11}^T A_{21} - A_{21}^T A_{11} = O \\
                A_{11}^T A_{22} - A^T_{21}A_{12} = I \\
                A^T_{12} A_{21} - A^T_{22} A_{11} = -I \\
                A^T_{12} A_{22} - A^T_{22} A_{12} = O.
            \end{cases}
        \]
    \end{itemize}
\end{lemma}
\begin{proof}
    Пусть
    \[
        M = \begin{pmatrix*} A_{11} & A_{12} \\ A_{21} & A_{22}
        \end{pmatrix*}
    \]
    где $A_{ij}$ -- блочные матрицы $m \times m$.
    Тогда симплектичность матрицы будет означать
    \[
        \begin{pmatrix*}
            A^T_{11} & A_{21}^T \\ \\ A^T_{12} & A_{22}^T
        \end{pmatrix*} \cdot \begin{pmatrix*} O & I \\ \\ -I & O
        \end{pmatrix*} \cdot \begin{pmatrix*} A_{11} & A_{12} \\ \\ A_{21} & A_{22}
        \end{pmatrix*} = \begin{pmatrix*} O & I \\ \\ -I & O
        \end{pmatrix*}.
    \]
    То есть
    \[
        \begin{pmatrix*}
            A^T_{11} & A_{21}^T \\ \\ A^T_{12} & A_{22}^T
        \end{pmatrix*} \cdot
        \begin{pmatrix*}
            A_{21} & A_{22} \\ \\ -A_{11} & -A_{12}
        \end{pmatrix*} =
        \begin{pmatrix*}
            O & I \\ \\ -I & O
        \end{pmatrix*}.
    \]
    Что эквивалентно следующей системе уравнений на матрицы:
    \[
        \begin{cases}
            A_{11}^T A_{21} - A_{21}^T A_{11} = O \\
            A_{11}^T A_{22} - A^T_{21}A_{12} = I \\
            A^T_{12} A_{21} - A^T_{22} A_{11} = -I \\
            A^T_{12} A_{22} - A^T_{22} A_{12} = O.
        \end{cases}
    \]
\end{proof}
Пусть у нас есть дифференциальная 1-форма следующего вида:
\[
    w = \bt{p}\cdot d\bt{q} + \tilde{\bt{q}}\cdot d\tilde{\bt{p}} = \bt{p}\cdot d\bt{q} + \tilde{\bt{q}} \cdot (\grad_{\bt{q}} \tilde{\bt{p}} d\bt{q} + \grad_{\bt{p}} \tilde{\bt{p}} d\bt{p}) = (\bt{p} + \tilde{\bt{q}}^T\grad_{\bt{q}} \tilde{\bt{p}})\cdot d\bt{q} + (\tilde{\bt{q}}^T\grad_{\bt{p}}\tilde{\bt{p}})d\bt{p}.
\]
(здесь $\tilde{\bt{p}}$ и $\tilde{\bt{q}}$ получены при помощи $\phi$). \\
Справедлива следующая теорема:
\begin{theorem}
    Для формы $w$ следующие условия эквивалентны:
    \begin{itemize}
        \item $w$ является полным дифференциалом некоторой функции $H$, то есть $w = dH$.
        \item $\phi$ -- симплектический диффеоморфизм.
    \end{itemize}
\end{theorem}
\begin{proof}
    1-форму $w$ можно записать как
    \[
        w = A \cdot d\bt{x}.
    \]
    где $A$ -- некоторый столбец. \\
    Заметим, что $w$ -- полный дифференциал некоторой функции $\Longleftrightarrow$ $\grad A$ -- симметрическая матрица (по крайней мере, в односвязной области: это следует из того, что в таком случае $dw = 0$ -- а всякая замкнутая форма в односвязной области точна).
    Ранее мы расписали какой вид имеет $w$:
    \[
        w = (\bt{p} + \tilde{\bt{q}}^T\grad_{\bt{q}} \tilde{\bt{p}})\cdot d\bt{q} + (\tilde{\bt{q}}^T\grad_{\bt{p}}\tilde{\bt{p}})d\bt{p}.
    \]
    То есть $A$ имеет вид
    \[
        A = \sum\limits_{j} \grad^T \tilde{p}_j \cdot \tilde{q}_j + \begin{pmatrix*} \bt{p} \\ 0
        \end{pmatrix*}.
    \]
    А значит её градиент равен:
    \[
        \grad A = \grad^T \tilde{\bt{p}} \grad \tilde{\bt{q}} + \sum\limits_{j} \dfrac{\partial^2 \tilde{p}_j}{\partial \bt{x} \partial \bt{x}} \tilde{q}_j + \begin{pmatrix*} O & I \\ O & O
        \end{pmatrix*}.
    \]
    Поскольку в природе не существует негладких функций, то второй компонент в градиенте симметричен. \\
    Рассмотрим матрицу
    \[
        B = \grad^T \tilde{\bt{p}} \grad \tilde{\bt{q}} + \begin{pmatrix*} O & I \\ O & O \end{pmatrix*}
    \]
    Покажем, что необходимым и достаточным условием симметричности $B$ является симплектичность $\phi$ (это завершит доказательство теоремы).
    Распишем дифференциал $\phi$:
    \[
        D\phi =
        \begin{pmatrix*}
        A_{11} & A_{12} \\ A_{21} & A_{22}
        \end{pmatrix*} =
        \begin{pmatrix*}
        D_{\bt{q}} \tilde{\bt{q}} & D_{\bt{p}} \tilde{\bt{q}} \\ D_{\bt{q}} \tilde{\bt{p}} & D_{\bt{p}} \tilde{\bt{p}}
        \end{pmatrix*}
    \]
    В то же время, в этих обозначениях, матрицу $B$ можно расписать как 
    \[
        B = 
        \begin{pmatrix*}
            D_{\bt{q}}^T \tilde{\bt{p}} \\ \\ D_{\bt{p}}^T \tilde{\bt{p}}
        \end{pmatrix*} \begin{pmatrix*} D_{\bt{q}} \tilde{\bt{q}} & D_{\bt{p}} \tilde{\bt{q}}
        \end{pmatrix*} + \begin{pmatrix*} O & I \\ \\ O & O
        \end{pmatrix*}.
    \]
    Нетрудно заметить, что условие
    \[
        B - B^T = O.
    \]
    в точности совпадает с требованиями предыдущей леммы, что и завершает доказательство теоремы.
\end{proof}
Рассмотрим функцию $S = S(\bt{q}, \tilde{\bt{p}})$ такую, что $\det \dfrac{\partial^2 S}{\partial \tilde{\bt{p}} \partial \bt{q}} \neq 0$.
\begin{theorem}[О производящей функции]
    Такая функция (называемая производящей) $S$ определяет преобразование $\phi: (\bt{q}, \bt{p}) \to (\tilde{\bt{q}}, \tilde{\bt{p}})$, где $\bt{p}$ и $\tilde{\bt{q}}$ определяются как
    \[
        \bt{p} = \partial_{\bt{q}} S, \quad \tilde{\bt{q}} = \partial_{\tilde{\bt{p}}} S,
    \]
    которое является симплектическим диффеоморфизмом. \\
    Более того, для всякого симплектического диффеоморфизма найдётся $S(\bt{q}, \tilde{\bt{p}})$ которая является производящей функцией для него.
\end{theorem}
\begin{proof} \leavevmode
    \begin{enumerate}
        \item При определённых таким образом $\bt{p}$ и $\bt{\tilde{q}}$ мы имеем
        \[
            w = \partial_{\bt{q}} S d\bt{q} + \partial_{\bt{\tilde{p}}} Sd\bt{\tilde{p}} = dS.
        \]
        А следовательно, по предыдущей лемме $\phi$, определяемый $S$, является симплектическим диффеоморфизмом.
        \item Если $\phi$ -- симплектический диффеоморфизм, то форма $w = dS$ и $S$ -- искомая производящая функция (опять-таки, подобный вывод делается исходя из предыдущей леммы).
    \end{enumerate}
\end{proof}
